\documentclass[11pt]{article}
\usepackage{amsmath,amssymb,amsthm}
\usepackage{fullpage}
\usepackage{enumitem}
\usepackage{algorithm}
\usepackage{algpseudocode}
\usepackage{optidef}
\usepackage{hyperref}
\usepackage{xcolor}

\title{A Memetic Genetic Algorithm for the Minimum Zero Forcing Set Problem}
\author{Thomas R. Cameron}
\date{\today}

\newtheorem{theorem}{Theorem}
\newtheorem{proposition}{Proposition}
\newtheorem{lemma}{Lemma}

\newcommand\abs[1]{\left|#1\right|}
\newcommand\ceil[1]{\left\lceil #1\right\rceil}
\newcommand\floor[1]{\left\lfloor #1\right\rfloor}
\newcommand\cl[1]{\operatorname{cl}\left(#1\right)}
\newcommand\shuffle[1]{\operatorname{shuffle}\left(#1\right)}


\begin{document}

\maketitle

\begin{abstract}
This report describes a memetic genetic algorithm for the minimum
zero forcing set problem. The algorithm combines evolutionary search
with a repair operator based on the zero forcing closure. Every
infeasible candidate solution encountered during the search is
converted into a feasible zero forcing set through a fort-guided
repair procedure. The minimal forts discovered during repair are
used to guide future crossover operations, generate generation-level
lower bounds, and produce structurally diverse solutions. Unlike
classical genetic algorithms, the proposed method continuously
extracts structural information from failed solutions and feeds this
information back into the evolutionary process.
\end{abstract}

\section{Introduction}

Let $G=(V,E)$ be a simple graph of order $n$. A set
$S \subseteq V$ is a \emph{zero forcing set} if repeated application
of the color change rule eventually colors every vertex of the graph.
The minimum cardinality of a zero forcing set is called the
\emph{zero forcing number} and is denoted by $Z(G)$.

The zero forcing problem is NP-hard, and therefore exact methods
may require substantial computational effort on large graphs.
Genetic algorithms provide a natural framework for exploring
large combinatorial search spaces. 
The objective of this report is to describe a specialized genetic
algorithm designed specifically for zero forcing. The algorithm
combines three key ideas:

\begin{enumerate}[label=(\roman*)]
\item every candidate solution is repaired into a feasible
zero forcing set,

\item failed candidate solutions generate minimal forts,

\item fort information is used to influence future generations.
\end{enumerate}

Because of the repair operator and local improvement procedures,
the method is more accurately described as a memetic algorithm
than a classical genetic algorithm.

\section{Preliminaries}
For a set $S \subseteq V$, let $\cl{S}$ denote the zero forcing closure of $S$.
If $\cl{S}=V$, then $S$ is a zero forcing set.
Otherwise, $F=V\setminus\cl{S}$ is a fort, that is, no vertex outside $F$ has exactly one neighbor in $F$.
The following characterization motivates the algorithm.

\begin{theorem}
A set $S \subseteq V$ is a zero forcing set if and only if
$S$ intersects every fort of $G$.
\end{theorem}

Thus every failed candidate solution provides information about
forts that must be intersected by any feasible solution.

\subsection{Solution Representation}
Each individual is represented by a subset $S \subseteq V$.
Internally, a binary encoding may be used, where
\[
x_v =
\begin{cases}
1 & v \in S,\\
0 & v \notin S.
\end{cases}
\]

\subsection{Vertex Weight}
The genetic algorithm will often take random samples of vertices to construct zero forcing sets.
When doing so, it is helpful to prioritize vertices that cover more forts.
Since there can be an exponential number of forts, we don't plan on storing all forts.
Instead, we will keep track of a vertex weight that represents the number of forts we have encountered.
Note that every time a failed zero forcing set is encountered, the complement of its closure is a fort. 
Moreover, in a single generation of the genetic algorithm we will encounter many forts; once the generation is complete, we update the vertex weight with the distinct forts encountered. 

Initially, each vertex $v\in V$ gets weight $w_{v}=0.01$.
Thus, the initial vertex weights correspond to a uniform distribution, and once forts are observed the corresponding weights change dramatically. 
Then, after each generation we take the collection of observed forts $\mathcal{F}$ and we update the vertex weights as follows
\[
w_{v} = \epsilon\cdot w_{v} + \sum_{F\in\mathcal{F}\colon v\in F}\frac{1}{\abs{F}},
\]
where $\epsilon=0.9$.
Note that the product $\epsilon\cdot w_{v}$ is used to prioritize recent generational fort information over older information.
Moreover, the term $\frac{1}{\abs{F}}$ prioritizes smaller forts over bigger forts. 

\subsection{Random Sampling Conventions}
Throughout the algorithm, we distinguish between uniform random choices and weight-biased random choices.

A random permutation of a set $A$, denoted by $\shuffle{A}$, is always taken uniformly over all permutations of $A$.
Similarly, a random sample of the population members is taken uniformly. 

By contrast, random samples of vertices from $V$, or from a specified subset $S\subseteq V$, may be taken using the current vertex weights $w_v$.
Unless explicitly stated otherwise, such vertex samples are weight-biased, with vertices of larger weight more likely to be selected.
In particular, for each $v\in S$, we select a uniform random number $\mu\in(0,1)$ and compute
\[
\mu^{1/a_{v}},
\]
where $a_{v}=\frac{w_{v}}{\sum_{u\in V}w_{u}}$ is the normalized weight of vertex $v$. 
Since $\mu\in(0,1)$ vertices with larger weights will correspond to bigger values.
A weight-biased random sample of size $k$ returns the $k$ vertices with largest values. 

\subsection{Population Initialization}
The algorithm maintains a population of size $p=10$.
Throughout the search, lower and upper bounds satisfy $lb \le Z(G) \le ub$.
Initially, the bounds may be taken as $lb=1$ and $ub=n$, although stronger bounds based on the graph structure are often available. 

To generate the population of size $p$, we define $\Delta p=\ceil{\frac{ub-lb}{p}}$.
Then, for each 
\[
k\in\{lb,lb+\Delta p,lb+2\Delta p,\ldots,lb+(p-1)\Delta p\}
\]
\begin{enumerate}
\item Generate a weight-biased random subset $S \subseteq V$ with
$|S|=k$.

\item Apply the repair procedure.

\item Insert the repaired solution into the population.
\end{enumerate}
Note that the repair procedure likely changes the size of the initial random sample by either adding vertices so that $S$ becomes a zero forcing set, or removing vertices so that $S$ is minimal.
Moreover, while no repaired zero forcing $S$ will be smaller than $lb$, it may be larger than $ub$; indeed, minimal zero forcing sets can be larger than the incumbent upper bound. 

\subsection{Fitness Scores}
Each population member gets a fitness score defines as follows
\[
f(S) = \alpha\cdot\left(\abs{V}-\abs{S}\right) + \frac{\beta}{\sum_{v\in V}w_{v}}\cdot\sum_{v\in S}w_{v},
\]
where $\alpha=1$ and $\beta=0.1$. 
Therefore, the fitness score first prioritizes the size of the zero forcing set and second the relative weight of the forts it hits. 

\subsection{Repair Procedure}
The repair procedure is an important aspect of this genetic algorithm because it both ensures that every member of the population is a zero forcing set and provides important fort information.
Given a candidate set $S$, compute its closure $C=\cl{S}$.
If $C=V$, then $S$ is already a zero forcing set and we apply the trim procedure to ensure $S$ is minimal.
Otherwise, $C$ is a failed zero forcing set and $F=V\setminus C$ is a fort.

Let $\shuffle{F}$ denote a uniformly random permutation of the vertices in $F$.
Also, let $z\in\{0,1\}$ denote a binary variable that indicates whether or not $S$ is a zero forcing set, initially $z=0$. 
Then, for each $u\in\shuffle{F}$
\begin{enumerate}
\item If $u\notin C$, then set $C' = C\cup\{u\}$. 

\item If $\cl{C'}= V$ and $z=0$, then update $S=S\cup\{u\}$ and $z=1$.

\item If $\cl{C'}\neq V$, then update $C=\cl{C'}$; furthermore, if $z=0$ then update $S=S\cup\{u\}$.
\end{enumerate}
Once all vertices in $F$ have been considered, the set $C$ is a maximal failed zero forcing set, $F=V\setminus C$ is a minimal fort, and $S$ is a zero forcing set.
We ensure that $S$ is a minimal zero forcing set by applying the trim procedure. 

\subsection{Trim Procedure}
The trim procedure takes a zero forcing set $S$ and removes as many vertices as possible while maintaining the zero forcing property. 
In particular, let $\shuffle{S}$ denote a uniformly random permutation of the vertices in $S$.
Then, for each $u\in\shuffle{S}$
\begin{enumerate}
\item Let $S'=S\setminus\{u\}$.

\item If $\cl{S'}=V$, then update $S=S'$. 
\end{enumerate}
Once all vertices in $S$ have been considered, the updated set $S$ is a minimal zero forcing set. 

\section{Generational Evolution}
The genetic algorithm evolves its population generation by generation. 
The initial population is considered generation $0$.
Let $g\geq 0$ denote the generation number and let 
\[
\mathcal{P}_{g} = \left\{S_{1}^{(g)},S_{2}^{(g)},\ldots,S_{10}^{(g)}\right\}
\]
denote the population at generation $g$, where $S_{1}^{(g)},S_{2}^{(g)},\ldots,S_{10}^{(g)}$ are the members of the population (zero forcing sets) ordered from largest to smallest fitness scores. 

The top $10\%$ of the population, $S_{1}^{(g)}$, are considered the elite members of the population at generation $g$. 
We update the remaining $90\%$ of the population through the selection, crossover, and replacement events. 

\subsection{Selection}
In the selection process, two members of the population are selected to construct a new member of the population via crossover.
The members are selected as follows: A uniformly random sample of size $3$ is taken from the population $\mathcal{P}_{g}$.
Then, all $\binom{3}{2}$ combinations are scored based on their combined fitness and symmetric difference:
\[
s\left(S_{i}^{(g)},S_{j}^{(g)}\right) = \frac{\gamma}{f(S_{1}^{(g)})}\cdot\left(f(S_{i}^{(g)})+f(S_{j}^{(g)})\right) + \delta\cdot\left(1 - \frac{\abs{S_{i}^{(g)}\cap S_{j}^{(g)}}}{\abs{S_{i}^{(g)}\cup S_{j}^{(g)}}}\right),
\]
where $\left(1-\frac{\abs{A\cap B}}{\abs{A\cup B}}\right)$ is the Jaccard distance between $A$ and $B$, $\gamma=1$, and $\delta=0.1$.

\subsection{Crossover}
Once two members of the population have been selected, say $S_{i}^{(g)}$ and $S_{j}^{(g)}$, the crossover process constructs a new member of the population.
Let $H\subseteq V\setminus\left(S_{i}^{(g)}\cup S_{j}^{(g)}\right)$ denote a weight-biased random sample of a random size.
Also, let $r=3$ and $\Delta r=\ceil{\frac{ub-lb}{r}}$.
Then, for each 
\[
k\in\left\{lb,lb+\Delta r,lb+2\Delta r,\ldots,lb+(r-1)\Delta r\right\}
\]
\begin{enumerate}
\item Generate a weight-biased random subset $S\subseteq S_{i}^{(g)}\cup S_{j}^{(g)}\cup H$ with $\abs{S}=k$.

\item Apply the repair procedure.
\end{enumerate}
After all $k$ have been considered, store the zero forcing set that has the highest fitness score. 

\subsection{Replacement}
Suppose that the selection and crossover processes have been repeated $8$ times.
Then, there are approximately $8r$ minimal forts that have been constructed from the repair function. 
Let $\mathcal{F}_{g}$ denote the (distinct) minimal forts that have been encountered during generation $g\geq 0$. 
We use these minimal forts to update the vertex weights and to construct one more zero forcing set to be added to the population.

Define $Z_{g}(G)$ to be the minimum size of a set of vertices $S$ such that $S\cap F\neq\emptyset$ for all $F\in\mathcal{F}_{G}$. 
Note that $Z_{g}(G)\leq Z(G)$.
Moreover, $Z_{g}(G)$ can be computed as a fort cover model IP over the forts in $\mathcal{F}_{g}$; since the number of minimal forts in $\mathcal{F}_{g}$ is at most $8r$, this model will likely finish quickly even on large graphs.
Once the model completes, we have an optimal set $S$ that hits every fort in $\mathcal{F}_{g}$.
While $S$ hits every fort in $\mathcal{F}_{g}$, it may not not be a zero forcing set of $G$. 
In this case, we apply the repair function to $S$.

At this point, we have constructed $9$ minimal zero forcing sets that are used to update population members $S_{2}^{(g)},\ldots,S_{10}^{(g)}$. 
Upon updating the members of the population, we are now in generation $g+1$.
At this point, we re-order the population from largest to smallest fitness, discard the forts observed in the previous generation, and we update $lb=\max\{lb,Z_{g}(G)\}$ and 
\[
ub = \min\left\{ub,\abs{S_{1}^{(g+1)}},\abs{S_{2}^{(g+1)}},\ldots,\abs{S_{10}^{(g+1)}}\right\}.
\]

\end{document}